\documentclass{article}
\usepackage[utf8]{inputenc}
\usepackage{a4wide}

\begin{document}

\section*{เซตย่อยต่อเนื่อง}

กำหนดอาร์เรย์เลขจำนวนเต็ม ( เริ่มต้นตำแหน่งที่ 0 )

ชื่อว่า\texttt{nums}ให้\textbf{หาจำนวนของเซตย่อยต่อเนื่อง}ของ\texttt{nums}โดย\ul{\textbf{เซตย่อยต่อเนื่องหมายถึงเซตย่อยที่สมาชิกทุกคู่ในนั้นมีค่าต่างกันไม่เกิน

2}}นั่นคือ



ให้\texttt{i},\texttt{i+1},...,\texttt{j}คือ{index ของเซต

โดยที่}\texttt{i\ \textless{}=\ i1,\ i2\ \textless{}=\ j},\texttt{0\ \textless{}=\ \textbar{}nums{[}i1{]}\ -\ nums{[}i2{]}\textbar{}\ \textless{}=\ 2}



\hfill\break



\textbf{ตัวอย่างเช่น}



\begin{quote}

\textbf{Input:}



\textbf{4}



\textbf{5 4 2 4}



\textbf{Output: 8}



\textbf{\ul{คำอธิบาย}}



เซตย่อยต่อเนื่องความยาว 1:\texttt{{[}5{]},\ {[}4{]},\ {[}2{]},\ {[}4{]}}\\



เซตย่อยต่อเนื่องความยาว 2:\texttt{{[}5,\ 4{]},\ {[}4,\ 2{]},\ {[}2,\ 4{]}}\\



เซตย่อยต่อเนื่องความยาว 3:\texttt{{[}4,\ 2,\ 4{]}}\\



เซตย่อยต่อเนื่องความยาว 4:\texttt{ไม่มี}



ดังนั้นจึงตอบ\texttt{4\ +\ 3\ +\ 1\ =\ 8}

\end{quote}



\textbf{\hfill\break

}



\subsection*{Input}
{บรรทัดแรกคือขนาดความกว้างของ{อาร์เรย์}}



{บรรทัดที่สองคืออาร์เรย์เลขจำนวนเต็มคั่นด้วยช่องว่าง}\\



\subsection*{Output}
{จำนวนของเซตย่อยต่อเนื่อง}\\



\end{document}
